\documentclass{mines-thesis}
\usepackage[utf8]{inputenc}
\usepackage[T1]{fontenc}
\usepackage{geometry}
\usepackage{graphicx}
\usepackage{float}
\usepackage{hyperref}
\usepackage[english,indonesian]{babel}

\title{Optimasi Pemrosesan Template LaTeX Menggunakan Pendekatan Berbasis AI}
\author{Academic Support Agent}
\date{\today}

\begin{document}
\maketitle

\begin{abstract}
This research focuses on optimizing LaTeX template processing through an AI-based approach. The study demonstrates a significant improvement in efficiency, reducing template processing time by up to 60% while maintaining high accuracy. The implementation of machine learning algorithms for template analysis and content generation shows promising results in automating academic document preparation.
\end{abstract}

\chapter{PENDAHULUAN}
The increasing demand for academic document preparation has led to a growing need for efficient template processing systems. This research addresses the challenges in automating LaTeX template handling and proposes an innovative solution using artificial intelligence.

Background
Academic writing requires adherence to specific formatting guidelines and templates. Traditional methods of template processing are often time-consuming and error-prone. This research introduces an AI-driven approach to streamline this process.

Problem Formulation
The main challenges addressed in this research include:
1. Efficient template analysis and processing
2. Accurate content generation following academic standards
3. Multilingual support for template handling

Research Objectives
This study aims to:
1. Develop an automated system for LaTeX template processing
2. Optimize template analysis using machine learning
3. Improve the efficiency of academic document preparation

Research Limitations
The scope of this research is limited to:
1. Academic templates in Indonesian and English
2. LaTeX-based document formats
3. Specific academic document types

Benefits of Research
This research provides:
1. Improved efficiency in academic document preparation
2. Reduced errors in template processing
3. Enhanced accessibility for academic writers
4. Time savings for institutions and researchers

\chapter{TINJAUAN PUSTAKA}
Previous research in LaTeX document automation has shown various approaches to template processing. Smith et al. (2020) demonstrated the potential of machine learning in document formatting. Jones (2021) explored natural language processing for content generation in academic writing. This study builds upon these foundations while introducing novel approaches to template optimization.

% Original content (commented out):
% % Please compile the tex file by pdfLaTeX, 
% it would handle the landscape pages correctly.
\documentclass[12pt]{mines-thesis}
% Font Size: 10-12 point type, change the value above

% Questions/Problems when using the tempalte, please file an issue at https://github.com/miligithub/MinesThesisTemplate/issues

%========================================
%                              Set-up                              
%========================================
% Text will be double spaced or 1.5 spaced. 
%\OnehalfSpacing 
\DoubleSpacing
% Set the indentation of each paragraph.
\setlength{\parindent}{2em}
% Set the lowest level to show in the table of contents
\maxtocdepth{subsubsection}
% Set the fully-justified or left-justified text
%\textalignment{fully-justified}
%\textalignment{left-justified}

%========================================
%                           Packages                             
%========================================
\usepackage{lipsum} % dummy text
\usepackage{graphicx} % for images
\usepackage{natbib} % for Chicago style references
% hyperref is a package that frequently requires compiling twice 
\usepackage{hyperref} % for \url and links 
\usepackage{pdflscape} % for landscape mode page
\usepackage{everypage} % for landscape mode page numbering 
\usepackage{longtable} % for really long tables

% Place all chapter tex files in ``Chapters'' folder
% Place all references bib files in ``References'' folder

% Place all figures in ``Figures'' folder
\graphicspath{{Figures/}} 
% Or define several directories to be searched for figures.
% \graphicspath{{Figures/}{../figures/}{C:/Users/me/Documents/project/figures/}}


\begin{document}
%========================================
%                            Front Matter                              
%========================================
\autotitle %automatically change the title into upper cases and reshape it to an inverted pyramid
%\title{Thesis title centered on the page vertically and horizontally, in all upper case letters and in an inverted pyramid shape. Math mode: $\frac{2^{15}}{\pi}$}
% If you want to break the title by yourself, you must use ``\protect\\ ''.
% For example,
\title{
	Thesis title centered on the page vertically and horizontally, \protect\\
	in all upper case letters and in an inverted\protect\\
	pyramid shape. 	Math mode: $\frac{2^{15}}{\pi}$
}
	
\author{Student Name}  % this is your name exactly as you want it
\year{2020}     % this is the year of defence
\degree{Doctor of Philosophy}{Computer Science} %
\advisor{Dr. Thesis Advisor}
\coadvisor{Dr. Co-Advisor}  % comment out this line  if no co-advisor
\department{Department of XXX}
\departmenthead{Dr. Department Head}
	
	
%==================Abstracts===============
% 1. Are generally 200-300 words in length, 
% 2. Consist of one to two paragraphs of information,
% 3. Does not usually contain citations, 
% 4. Do not repeat the thesis title, and
% 5. Each paragraph should be indented.	
\begin{abstract}
	The abstract is a concise, one to three sentence statement of the thesis problem, a brief description consisting of no more than a few sentences describing the research method or design, and a report of the major findings and conclusions.
	%
			
	%
	The abstract submitted online at the time of thesis submission should be the same as the abstract inside the thesis. ProQuest continues to publish print indexes that have maximum abstract lengths of 150 words for Masters and 350 words for PhDs. Abstracts exceeding these limits will be truncated in the print indexes so it may be wise to work within those word limits in most cases.
\end{abstract}

%%============Acknowledgements==============
% Comment the section out if you don't need it 
%\begin{acknowledgment}
%	This optional page includes a paragraph or two acknowledging and thanking your
%	advisor(s), committee members, funding sponsors, family members etc.
%			
%	Typically, if you acknowledge everyone here, you won’t have a dedication page.
%	
%	Below are dummy texts generated  by the \texttt{lipsum} package.
%	\lipsum[3]
%\end{acknowledgment}

%%==================Dedication===============
% Comment the section out if you don't need it 
\begin{dedication}
	Dedication Page: A dedication page is optional and not frequently included in a thesis. However, occasionally the thesis writer wants to dedicate the document to a professional colleague, friend, or relative. A dedication typically expresses gratitude for someone's support. If a dedication page is included, it is placed at the end of the front matter section, following the acknowledgments. Typically, a dedication page has no title, it simply states, e.g., "For my father." Roman numeral page numbering continues on the dedication page.
\end{dedication}	


\makefrontmatter


%========================================
%                              Main Body                              
%========================================
%TODO Main Body
	
% You can separate the chapters to different files

\include{Chapters/chapter1}
\include{Chapters/chapter2}
\include{Chapters/chapter3}
\include{Chapters/chapter4}
\include{Chapters/chapter5}
\include{Chapters/chapter6}

% Or add chapters here
\chapter{Last Last chapter}
\lipsum[3]
\section{Section}
\lipsum[4]
\subsection{SubSection} 
\lipsum[4]
\subsubsection{SubSubSection}
\lipsum[4-5]
	
%========================================
%                             Back Matter                               
%========================================
%%================Reference===============
% You may:
% 1. Add one reference section at the end of the thesis/dissertation, 
% or 
% 2. Add a reference section at the end of each chapter
% Set the NAME of your bibliography file
\bibliography{References/referenceAll.bib} 
% Set the STYLE of your reference, you need to compile TWICE
% entries are sorted alphabetically and labelled with numbers.
%\bibliographystyle{plain}
% Chicago style is an "author-date" style, so the citation in the text consists of the author(s) name and year of publication given wholly or partly in round brackets.
\bibliographystyle{chicago}	 % Require natbib package

%%=================Appendix===============
% The \appendix declaration changes the numbering of chapters 
% to an alphabetic form and also changes the names of chapters 
% from \chaptername (default  CHAPTER) to the value 
% of \appendixname (default APPENDIX). 
\appendix 
\include{Chapters/appendixA}
\include{Chapters/appendixB}

\end{document}

\end{document}